% =====================================================================
% Global mapping of concentrated animal feeding operations
% Publication-ready draft. Compile with:
%     pdflatex main && bibtex main && pdflatex main && pdflatex main
% Bibliography: references.bib
% =====================================================================
\documentclass[11pt,a4paper]{article}

\usepackage[utf8]{inputenc}
\usepackage[T1]{fontenc}
\usepackage{lmodern}
\usepackage[margin=1in]{geometry}
\usepackage{amsmath,amssymb}
\usepackage{booktabs}
\usepackage{multirow}
\usepackage{graphicx}
\usepackage{siunitx}
\usepackage{longtable}
\usepackage{listings}
\usepackage{xcolor}
\usepackage[round,authoryear]{natbib}
\usepackage[colorlinks=true,allcolors=blue]{hyperref}
\usepackage{microtype}

\sisetup{detect-all,group-minimum-digits=4}

\lstset{
  basicstyle=\ttfamily\footnotesize,
  breaklines=true,
  frame=single,
  framesep=4pt,
  columns=fullflexible,
  keepspaces=true,
  showstringspaces=false
}

\newcommand{\code}[1]{\texttt{\small #1}}
% Marks a design choice that is currently an inherited default rather than
% an ablated decision. Every occurrence maps to an entry in the companion
% experiment plan (paper/experiments_justification_plan.md).
\newcommand{\unablated}[1]{\textsuperscript{\dag}\,(#1)}

\title{Global Mapping of Concentrated Animal Feeding Operations from
       Sentinel-2 Imagery: Pipeline, Ablations, and Measurement Validity}

% TODO(author): confirm full author list, affiliations, and corresponding address.
\author{Filip Zawadka \\ \small Rachel~[surname TBC]}
\date{\today}

\begin{document}
\maketitle

\begin{abstract}
Concentrated animal feeding operations (CAFOs) are a significant source of
regional water and air pollution and a focal point for zoonotic-disease
surveillance, yet no openly available inventory covers them at global scale.
We describe an end-to-end pipeline that detects and types livestock facilities
worldwide from freely available Sentinel-2 imagery. The pipeline is two-stage:
a geometric stage clusters candidate barn-like structures from the Google Open
Buildings footprint database using morphometric filters, and a convolutional
stage classifies a 1.28\,km Sentinel-2 patch centred on each candidate cluster
into \textsc{NotFarm}, \textsc{Poultry}, \textsc{Pigs}, or \textsc{Cattle}. The
production model is a ResNet-50 initialised from SoftCon, a Sentinel-2--native
self-supervised backbone, adapted to a nine-channel input (six surface-reflectance
bands plus three spectral indices) by band-name--matched first-convolution
transfer. We score \num{157099} of \num{157102} candidate clusters across 167
countries.

Beyond the system description, we report a systematic experimental record
covering backbone choice, spatial context, class-imbalance interventions, label
acquisition strategy, and regularisation. Three findings recur. First, after
de-duplication and bootstrap confidence intervals, the only comparison that
separates robustly from noise is Sentinel-2--native self-supervised pretraining
versus ImageNet initialisation ($\approx\!+0.10$ macro-F1, roughly twice the
confidence-interval half-width); a nine-channel ImageNet control shows the gain
is attributable to pretraining rather than band count, and the six leading
Sentinel-2--native configurations are statistically tied. Second, class-balanced
sampling is a reproducible in-domain/out-of-domain trade rather than an
improvement, gaining on held-out data from training countries while losing on
unseen countries. Third, adding negative labels alone shifts the decision prior
without improving ranking quality: area under the ROC curve is unchanged while
both false-positive rate and recall fall together.

We also document, rather than omit, two data-integrity problems that materially
affected earlier results: a patch-store key misalignment that paired
approximately 34.6\% of labelled candidates with imagery from the wrong location,
and a systematic label-source shift between training and evaluation data. We
report corrected numbers, show that two previously published verdicts reverse
under clean data, and describe the invariants now enforced in code. Finally, we
argue that the dominant bottleneck in this problem is measurement rather than
modelling. Training the same configuration under five random seeds gives a
standard deviation of $0.033$ macro-F1 on the standard evaluation slice---a
decision band wider than nearly every intervention previously ranked on it---and
we show that roughly 80\% of that variance comes from a single class with six
held-out examples. Re-running twenty single-lever ablations against a purpose-built
blind benchmark, on which the same model is reproducible to $\sigma = 0.0008$,
we find that exactly one inherited default is worth changing, that a third of the
input channels contribute nothing measurable, and that two settings previously
believed to matter do not.
\end{abstract}

\noindent\textbf{Keywords:} concentrated animal feeding operations; Sentinel-2;
self-supervised pretraining; remote sensing; class imbalance; spatial
cross-validation; measurement validity.

\vspace{1em}
\noindent\rule{\linewidth}{0.4pt}
\noindent\footnotesize
Design choices marked \unablated{E$x$.$y$} are inherited defaults that are
\emph{not} yet supported by a controlled ablation in this work; each identifier
refers to the corresponding entry in the companion experiment plan
(\code{paper/experiments\_justification\_plan.md}).
\normalsize
\noindent\rule{\linewidth}{0.4pt}

% =====================================================================
\section{Introduction}
% =====================================================================

Intensive livestock production is concentrated in facilities that are large,
spatially clustered, and visually distinctive from above: long, narrow,
regularly spaced barns, often accompanied by waste lagoons. These facilities
are of standing interest to environmental regulators, because manure lagoons
and land application drive nutrient loading in surface and ground water; to
public-health agencies, because high-density confinement is a recognised
interface for zoonotic transmission; and to animal-welfare and food-system
researchers, who need facility counts to estimate populations under
confinement. In several jurisdictions the necessary inventory does not exist,
is not public, or is maintained only as an address list without geometry.

Remote sensing is an attractive substitute, and prior work has demonstrated it
convincingly at national scale. \citet{handannader2019cafo} detect CAFOs in
North Carolina from aerial imagery and show that model-assisted enforcement
targeting substantially outperforms complaint-driven inspection.
\citet{robinson2022poultry} map industrial poultry operations across the
contiguous United States using high-resolution aerial imagery. Both rely on
imagery products---sub-metre aerial orthophotography such as the USDA National
Agriculture Imagery Program---that have no global equivalent. Extending this
class of system worldwide therefore requires working from a sensor with global,
repeated, free coverage. Sentinel-2 \citep{drusch2012sentinel2} provides
10\,m multispectral imagery with a five-day revisit, which is roughly two
orders of magnitude coarser in area per pixel than the aerial imagery used in
prior CAFO work. At 10\,m, an individual poultry barn is on the order of
$2 \times 15$ pixels: the object is resolvable as an oriented bright
rectangle, but its texture, roof detail, and ventilation hardware are not.

This paper describes a pipeline built around that constraint, and reports what
we learned attempting to make it work at global scale. Our contributions are:

\begin{enumerate}
  \item \textbf{A two-stage global pipeline.} Rather than searching imagery
  exhaustively, we first generate candidates geometrically from the Google Open
  Buildings footprint database \citep{sirko2021openbuildings} using
  morphometric filters calibrated on a clean reference region, then classify a
  Sentinel-2 patch centred on each candidate cluster. This converts a global
  dense-search problem into a \num{157102}-candidate classification problem
  spanning 167 countries (Section~\ref{sec:data}).

  \item \textbf{A four-class facility-typing model.} We adapt a
  Sentinel-2--native self-supervised ResNet-50 \citep{wang2024softcon,he2016resnet}
  to a nine-channel input via band-name--matched first-convolution transfer, and
  train it to distinguish \textsc{NotFarm}, \textsc{Poultry}, \textsc{Pigs}, and
  \textsc{Cattle} (Section~\ref{sec:methods}).

  \item \textbf{A systematic ablation record with honest uncertainty.} We report
  a controlled comparison of backbones, spatial context, band composition,
  optimisation, augmentation, imbalance interventions, label-acquisition
  strategies, and regularisation. Crucially, we calibrate every comparison
  against measured run-to-run variance rather than against zero, and report
  which effects survive. Most do not (Section~\ref{sec:results}).

  \item \textbf{A data-integrity and measurement-validity methodology.} We
  document a patch-store key misalignment that silently contaminated
  approximately 34.6\% of labelled candidates for several months, quantify its
  effect, show that it \emph{reversed} two experimental verdicts, and describe
  the code-level invariants that now prevent it. We also quantify a
  train/evaluation label-source shift and the resulting ceiling on measurable
  progress (Section~\ref{sec:validity}).
\end{enumerate}

A theme runs through the experimental sections and we state it at the outset:
for most of this project's history, the measurement apparatus has been less
precise than the effects being measured. At an evaluation-set size of
$n=523$, the 95\% confidence interval on macro-F1 is approximately $\pm0.05$,
while the interventions under test---a different sampler, a different crop
size, a regularisation bundle---move the metric by $0.01$ to $0.05$. Reporting
such comparisons as wins would be unjustified. We instead report them as ties,
identify the single comparison that does separate, and treat improving the
benchmark as the highest-priority experiment rather than a housekeeping task.

% =====================================================================
\section{Related work}
% =====================================================================

\paragraph{CAFO and livestock-facility detection.}
\citet{handannader2019cafo} established the template for this task: detect
facilities with a convolutional network over aerial imagery, and evaluate not
only classification metrics but downstream regulatory utility.
\citet{robinson2022poultry} scaled poultry-barn detection across the
contiguous United States, emphasising the operational cost of false positives
in a screening deployment. Both operate at sub-metre resolution in a single
country with a reliable registry for supervision. Our work differs on all three
axes: 10\,m imagery, 167 countries, and heterogeneous supervision that mixes
national registries with visual annotation. The resolution change in particular
means our model cannot rely on the fine texture cues available to prior work,
and our results should be read as a different operating point---broad,
low-cost, globally uniform screening---rather than as a direct comparison.

\paragraph{Self-supervised pretraining for Earth observation.}
ImageNet initialisation transfers poorly to multispectral satellite input,
both because the domain gap is large and because non-RGB bands have no
pretrained counterpart. A line of work pretrains directly on Sentinel imagery:
SSL4EO-S12 \citep{wang2023ssl4eo} applies momentum contrast
\citep{he2020moco} to a large multitemporal Sentinel-1/2 corpus, and SoftCon
\citep{wang2024softcon} adds multi-label--guided soft contrastive learning,
reporting stronger performance at equal backbone capacity. TorchGeo
\citep{stewart2022torchgeo} distributes these weights with band metadata, which
our first-convolution adaptation depends on. GEO-Bench
\citep{lacoste2023geobench} provides the standardised comparison we used to
shortlist candidate backbones, alongside broader multimodal foundation models
such as DOFA \citep{xiong2024dofa} and CROMA \citep{fuller2023croma}. Common
downstream benchmarks in this literature---EuroSAT \citep{helber2019eurosat},
BigEarthNet \citep{sumbul2019bigearthnet}---are scene-classification tasks with
large, balanced label sets, which is worth noting: our task is an object-centred,
severely imbalanced classification with a long tail, and we find the benchmark
ranking only partially predicts our downstream ranking.

\paragraph{Long-tailed recognition.}
Our label distribution is severely skewed (roughly 50:1 \textsc{Poultry} to
\textsc{Cattle} in training). We evaluate the standard interventions:
class-weighted and class-balanced sampling; classifier retraining with a frozen
representation (cRT) \citep{kang2020decoupling}; train-time logit adjustment
and its post-hoc variant \citep{menon2021logit}; and focal loss
\citep{lin2017focal}. A recurring result in that literature---that decoupling
representation learning from classifier balancing outperforms end-to-end
resampling---holds only partially here, and we report the conditions under
which it flipped sign.

\paragraph{Regularisation.}
We tested label smoothing \citep{szegedy2016rethinking} together with cutout
\citep{devries2017cutout} and a longer schedule. \citet{muller2019labelsmoothing}
note that label smoothing improves accuracy and calibration in many settings but
can distort the confidence structure of the penultimate representation; our
measurement was that it worsened expected calibration error on the slice we care
about, which we report in Section~\ref{sec:v10}. Optimisation follows decoupled
weight decay \citep{loshchilov2019adamw} with cosine annealing
\citep{loshchilov2017sgdr}.

\paragraph{Evaluation validity in spatial problems.}
Random train/test splits over spatially autocorrelated data inflate apparent
performance, because held-out points are often near training points
\citep{roberts2017spatialcv,ploton2020spatial}. This applies directly to our
setting, where candidate clusters are dense in production regions; we quantify
the exposure in Section~\ref{sec:measurement}. We report bootstrap confidence
intervals \citep{efron1987bootstrap}, expected calibration error
\citep{guo2017calibration}, and Matthews correlation coefficient
\citep{chicco2020mcc} alongside F1, the last because our slices vary widely in
class balance.

% =====================================================================
\section{Data}
\label{sec:data}
% =====================================================================

\subsection{Candidate generation}
\label{sec:candidates}

The unit of analysis is a \emph{building cluster}, not a pixel or a tile. This
choice is what makes global coverage tractable: instead of running a detector
over all land area, we enumerate structures that are already known to exist and
filter them to those that are barn-shaped.

We start from Google Open Buildings \citep{sirko2021openbuildings}, a
continental-scale footprint database derived from high-resolution satellite
imagery, tiled by H3 hexagons at resolution~3. Footprints in urban areas are
removed. For each remaining building we compute morphometric descriptors---area,
length, width, aspect ratio, orientation---and apply the filters in
Table~\ref{tab:morphometry}. Thresholds were calibrated on the Delmarva
peninsula (``DMV''), a dense, well-documented United States broiler-production
region that serves throughout this project as the clean reference distribution.

\begin{table}[t]
\centering
\small
\caption{Morphometric candidate-selection parameters. Values are the operative
thresholds in the candidate-generation stage. The parameter set was calibrated
on the Delmarva reference region and applied globally without per-country
tuning.}
\label{tab:morphometry}
\begin{tabular}{llr}
\toprule
Parameter & Meaning & Value \\
\midrule
\code{MIN\_SIZE} / \code{MAX\_SIZE} & building footprint area & \SI{850}{}--\SI{5200}{\metre\squared} \\
\code{MIN\_AR} / \code{MAX\_AR}     & aspect ratio (elongation) & 4.5--15 \\
\code{MIN\_WIDTH} / \code{MAX\_WIDTH} & building width & \SI{10}{}--\SI{28}{\metre} \\
\code{MAX\_LENGTH}                  & building length & \SI{205}{\metre} \\
\code{BLDG\_SEP}                    & edge-linkage clustering distance & \SI{50}{\metre} \\
\code{CLUSTER\_MIN}                 & minimum buildings per cluster & 2 \\
\code{TOTAL\_AREA}                  & minimum total cluster footprint & \SI{2500}{\metre\squared} \\
\code{DIST\_TO\_WATER}              & lagoon-proximity search radius & \SI{150}{\metre} \\
\code{WATER\_PROBABILITY}           & Dynamic World water threshold & 0.1 \\
\code{FARM\_TO\_CLUSTER}            & registry-point to cluster match radius & \SI{100}{\metre} \\
\bottomrule
\end{tabular}
\end{table}

Buildings whose edges lie within \SI{50}{\metre} of one another are linked into
a single cluster; clusters of at least two buildings and at least
\SI{2500}{\metre\squared} total footprint become candidates. A lagoon feature is
derived from Dynamic World \citep{brown2022dynamicworld} water probability
within \SI{150}{\metre} of the cluster. Registry points are attached to clusters
within \SI{100}{\metre}.

This stage is deliberately high-recall and low-precision with respect to the
final task: it produces \num{157102} candidate clusters worldwide, of which
roughly three quarters of the labelled subset are in fact farms. That base rate
is important context for every metric in this paper, and we state it explicitly
in Section~\ref{sec:if}: a trivial classifier that labels every candidate as a
farm achieves precision $0.767$ at recall $1.0$. The convolutional stage must be
evaluated against that baseline, not against a 50\% prior.

\subsection{Labels}
\label{sec:labels}

Labels come from two qualitatively different sources, and the distinction
matters enough that we return to it in Section~\ref{sec:labelshift}.

\emph{Registry labels} are derived from official or curated facility lists:
United States state registries (including the Delmarva and Iowa sets), Mexican
and Chilean official registries, Thailand and Brazil facility data. A registry
point within \SI{100}{\metre} of a candidate cluster labels that cluster with
the registry's species, subject to animal-count thresholds (\num{20000} birds
for poultry, 400 head for pigs) to exclude smallholdings.

\emph{Visual labels} are assigned by human annotators inspecting high-resolution
basemap imagery for a candidate, and include explicitly adjudicated hard
negatives---greenhouses, warehouses, industrial sheds, plastic-covered
horticulture---that the morphometric stage cannot reject.

The upstream taxonomy is finer than the model's (it distinguishes broiler, egg,
and unspecified poultry, and carries \textsc{Mixed}, \textsc{Other},
\textsc{Unknown}, \textsc{PigsOrPoultry}, and \textsc{Ambiguous} categories). We
map it to four classes: \textsc{NotFarm}, \textsc{Poultry} (all poultry
subtypes), \textsc{Pigs}, and \textsc{Cattle}. Rows whose upstream label is
ambiguous or non-specific are dropped from supervised training but retained with
label $-1$ so that they are still scored at inference. Table~\ref{tab:labels}
gives the composition of the current label release.

\begin{table}[t]
\centering
\small
\caption{Label composition of the current release (\code{all\_clusters\_v9}).
\num{157102} clusters across 167 countries, of which \num{32214} carry a label.}
\label{tab:labels}
\begin{tabular}{lrlr}
\toprule
Upstream label & $n$ & Upstream label & $n$ \\
\midrule
NotFarm                  & \num{13778} & Poultry: Eggs      & \num{774}  \\
Poultry: Unspecified     & \num{8005}  & Cattle             & \num{228}  \\
Poultry: Meat Chickens   & \num{4265}  & Ambiguous          & \num{182}  \\
Pigs                     & \num{2576}  & PigsOrPoultry      & \num{111}  \\
Farm: Unknown            & \num{2172}  & Mixed              & \num{103}  \\
                         &             & Other              & \num{20}   \\
\bottomrule
\end{tabular}
\end{table}

The label set was built in rounds, and the rounds are themselves an experimental
variable (Section~\ref{sec:labelrounds}). Round~2 added \num{962} adjudicated
\textsc{NotFarm} corrections across 106 countries---exclusively negatives.
Round~3 (in two stages) promoted farm positives from the qualitative-evaluation
pool into training under a per-country cap, first at 50 and then at 70 per
country, with an 80:20 train/validation assignment.

\subsection{Splits}
\label{sec:splits}

Split assignment is explicit and upstream: each cluster carries a
\code{cnn\_split\_assigned} value, and the training code treats it as the single
source of truth. No random split is computed, and no country whitelist or
region-balancing logic is applied on this path. This replaced an earlier
random-number-generator--based splitter whose behaviour changed between runs and
which, at one point, silently dropped an entire class (Section~\ref{sec:validity}).

The splits carry distinct semantics rather than being three draws from one pool:

\begin{description}\itemsep2pt
  \item[\textnormal{\textsc{train} / \textsc{val} / \textsc{test}}] Standard
  fitting, model selection, and held-out measurement. \textsc{test} is drawn
  from the same five heavily-labelled countries as \textsc{train}, so it is the
  most optimistic slice.
  \item[\textsc{eval}] Held-out clusters from training countries, weighted
  toward visually adjudicated cases. Measures in-domain generalisation to harder,
  differently-sourced labels.
  \item[\textsc{generalization}] Countries with \emph{zero} training rows
  (Bangladesh, Albania, Nigeria, DR Congo). The out-of-domain slice, and the one
  closest to deployment conditions.
  \item[\textsc{qual\_eval}] A large adjudicated pool used for false-positive
  and recall estimation at scale. Comparisons across model versions use
  \textsc{qual\_eval\_common}, the intersection of the pool across all versions,
  which no model in the comparison trained on.
  \item[\textsc{predict}] Unlabelled candidates---the deployment target.
\end{description}

\begin{table}[t]
\centering
\small
\caption{Split composition of the current release. \textsc{Pigs} is absent from
\textsc{generalization} entirely, and \textsc{Cattle} support is in the single
digits on every held-out slice; both facts bound what can be claimed about those
classes.}
\label{tab:splits}
\begin{tabular}{lrrrrrr}
\toprule
Split & NotFarm & Poultry & Pigs & Cattle & Total & Countries \\
\midrule
\textsc{train}          & \num{2750} & \num{7616} & \num{1576} & 153 & \num{12095} & 139 \\
\textsc{val}            & 621        & \num{1662} & 356        & 33  & \num{2672}  & 113 \\
\textsc{test}           & 425        & \num{1344} & 300        & 25  & \num{2094}  & 5   \\
\textsc{eval}           & 135        & 319        & 62         & 7   & 523         & 5   \\
\textsc{generalization} & 171        & 251        & 0          & 4   & 426         & 4   \\
\textsc{qual\_eval}     & \num{9676} & \num{1852} & 282        & 6   & \num{11816} & 132 \\
\textsc{predict}        & \multicolumn{4}{c}{unlabelled} & \num{124888} & 167 \\
\bottomrule
\end{tabular}
\end{table}

Training data are geographically concentrated: the United States contributes
\num{6760} training rows, Mexico \num{1614}, Thailand 575, Brazil 454, and Chile
330, with a long tail of roughly 134 further countries contributing small
numbers. This concentration is why country-level rebalancing is disabled
(Section~\ref{sec:imbalance}) and why the \textsc{generalization} slice is the
metric we weight most heavily in deployment decisions.

A parallel column, \code{if\_split\_assigned}, defines splits for the geometric
Isolation Forest baseline (Section~\ref{sec:if}); the convolutional pipeline
passes it through to outputs but never trains on it.

\subsection{Imagery}
\label{sec:imagery}

For each candidate we extract a \(128 \times 128\) pixel patch at \SI{10}{\metre}
resolution---a \SI{1.28}{\kilo\metre} square---centred on the cluster centroid.
The patch grid is constructed in the local UTM zone with an explicit affine
transform, and pixels are requested directly from Google Earth Engine
\citep{gorelick2017gee} as arrays.

The source collection is \code{COPERNICUS/S2\_SR\_HARMONIZED}, harmonised
Level-2A surface reflectance \citep{mainknorn2017sen2cor}. We take an
\emph{annual median composite} over calendar year 2023, filtering scenes to
\code{CLOUDY\_PIXEL\_PERCENTAGE} below 15\%\unablated{E1.3}. Six bands are
retained---B2, B3, B4 (visible), B8 (NIR), B11, B12 (SWIR)---and three
normalised-difference indices are computed server-side:
\[
  \mathrm{NDVI}=\frac{B8-B4}{B8+B4},\quad
  \mathrm{NDBI}=\frac{B11-B8}{B11+B8},\quad
  \mathrm{NDWI}=\frac{B3-B8}{B3+B8},
\]
giving a nine-channel \code{float32} tensor per candidate\unablated{E1.1}. NDVI
separates barn roofs from surrounding vegetation, NDBI responds to built
surfaces, and NDWI is intended to expose the manure-lagoon signature that also
motivates the geometric water feature.

The annual median is a deliberate trade. It suppresses cloud, haze, and
transient agricultural change, and yields one stable image per site; it also
discards all phenological and temporal signal, which we expect is part of the
ceiling on distinguishing facility types (Section~\ref{sec:limitations}). A
per-pixel scene-classification cloud mask and a least-cloudy compositing mode
are implemented but were never evaluated: the one run configured to test them
was lost to a storage-quota failure and never repeated\unablated{E1.3}.

Patches are cached on disk keyed by candidate identifier and by an
\emph{imagery configuration hash}---a digest over bands, indices, compositing
mode, cloud threshold, and date range. The hash is both the cache key and the
inference-time filter, ensuring a model is never scored on patches built under a
different imagery configuration.

\paragraph{Cloud-blocked candidates and imagery tiers.}
Under the primary configuration, \num{2197} candidates could not be imaged at
all: in chronically cloudy tropical regions (Ecuador, Malaysia, Indonesia,
Brazil, the Philippines, Colombia, Panama, Costa Rica), \emph{no} Sentinel-2
scene in calendar 2023 passes a 15\% scene-cloud filter, so the composite is
empty. Retrying under the identical configuration recovered three. We therefore
relaxed the imagery configuration in two stages and recorded the provenance of
each recovered patch in an \code{imagery\_tier} field
(Table~\ref{tab:tiers}). Widening the window to 2022--2024 while holding the
cloud threshold at 15\% recovered \num{1810} candidates at full quality; relaxing
the threshold to 60\% recovered the remaining 384. Predictions on the
relaxed-cloud tier skew heavily toward \textsc{Poultry} (approximately 79\%),
consistent with the out-of-distribution bias discussed in
Section~\ref{sec:deployment}, and are flagged as screening-quality rather than
treated as equivalent to the primary tier.

\begin{table}[t]
\centering
\small
\caption{Imagery tiers in the full-world release. Coverage is
\num{157099}/\num{157102} candidates (99.998\%).}
\label{tab:tiers}
\begin{tabular}{llr}
\toprule
Tier & Imagery configuration & Candidates \\
\midrule
\code{primary\_2023}      & 2023 median, scene cloud $<15\%$      & \num{154905} \\
\code{clean\_2022\_2024}  & 2022--2024 median, scene cloud $<15\%$ & \num{1810}   \\
\code{relaxed\_cloud60}   & 2022--2024 median, scene cloud $<60\%$ & 384          \\
\midrule
\multicolumn{2}{l}{\textit{Total scored}} & \num{157099} \\
\bottomrule
\end{tabular}
\end{table}

% =====================================================================
\section{Methods}
\label{sec:methods}
% =====================================================================

\subsection{Architecture and channel adaptation}
\label{sec:arch}

The production model is a ResNet-50 \citep{he2016resnet} initialised from
SoftCon \citep{wang2024softcon}, a Sentinel-2--native backbone pretrained with
multi-label--guided soft contrastive learning and distributed through TorchGeo
\citep{stewart2022torchgeo}. We also evaluate SSL4EO-S12 MoCo weights
\citep{wang2023ssl4eo,he2020moco} and ImageNet initialisation; results are in
Section~\ref{sec:backbone}. Backbone selection was shortlisted using GEO-Bench
linear-probe results \citep{lacoste2023geobench}, on which SoftCon leads
comparable backbones at 64-pixel input\unablated{E1.6}.

Adapting a pretrained backbone to nine channels is not a matter of widening the
first convolution. The pretrained weights are ordered by the Sentinel-2 13-band
Level-1C convention, \([B1, B2, B3, B4, B5, B6, B7, B8, B8A, B9, B10, B11, B12]\),
whereas our input is \([B2, B3, B4, B8, B11, B12, \mathrm{NDVI}, \mathrm{NDBI},
\mathrm{NDWI}]\). Copying weights positionally would place the pretrained
\(B1\) filter on our \(B2\) input, and so on---a silent misalignment across every
channel. We therefore adapt by \emph{band name}: each input channel that has a
pretrained counterpart receives that counterpart's filter weights, and the three
index channels, which have none, are initialised to the mean of the pretrained
first-layer weights. The classification head is replaced with a fresh linear
layer over the 2048-dimensional pooled feature.

\subsection{Preprocessing and augmentation}

Each sample is loaded, non-finite values are zero-filled, spectral bands are
scaled by \(1/10000\) to surface reflectance, augmentation is applied, the patch
is centre-cropped to 64 pixels (\SI{640}{\metre})\unablated{E1.2}, the channel
subset is taken, and per-channel standardisation is applied.

Two normalisation regimes exist because the backbones demand different ones.
SSL4EO MoCo weights expect raw scaled reflectance with no standardisation;
SoftCon expects per-band \(z\)-scoring, which is its official transform. Per-channel
statistics are computed once over 512 randomly drawn training patches with
augmentation disabled, persisted to disk, and reloaded at inference---so that a
scoring shard uses the statistics of the run that trained the checkpoint, never
its own.

The augmentation recipe (Table~\ref{tab:augs}) is applied only to training
samples. It was assembled incrementally and has never been ablated
term-by-term\unablated{E1.5}. One inconsistency is worth flagging: per-band gain
jitter perturbs the spectral bands, but the index channels are \emph{not}
recomputed afterwards, so NDVI, NDBI, and NDWI are mildly inconsistent with the
bands they were derived from in roughly 30\% of training samples. A
\code{recompute\_indices} option exists and is disabled.

\begin{table}[t]
\centering
\small
\caption{Training augmentation recipe. Cutout is disabled in the production
model; it was enabled only in the rejected regularisation bundle
(Section~\ref{sec:v10}). No term in this table has been ablated individually.}
\label{tab:augs}
\begin{tabular}{llll}
\toprule
Augmentation & Probability & Parameters & Applies to \\
\midrule
Horizontal flip        & 0.5  & ---                       & all channels \\
Vertical flip          & 0.5  & ---                       & all channels \\
Rotation $90^\circ k$  & 0.75 & $k \in \{1,2,3\}$         & all channels \\
Continuous rotation    & 0.3  & $\pm 15^\circ$, reflect   & all channels \\
Random resized crop    & 0.3  & scale $0.8$--$1.0$        & all channels \\
Brightness jitter      & 0.5  & $\times[0.85, 1.15]$      & spectral only \\
Per-band gain jitter   & 0.3  & $\times[0.95, 1.05]$      & spectral only \\
Gaussian noise         & 0.3  & $\sigma = 0.02$           & spectral only \\
Channel dropout        & 0.1  & $\leq 1$ channel zeroed   & spectral only \\
Cutout \emph{(off)}    & ---  & $2 \times 16$\,px holes   & all channels \\
\bottomrule
\end{tabular}
\end{table}

\subsection{Training}
\label{sec:training}

The canonical recipe, held byte-identical across the label-round series so that
data effects are not confounded with recipe effects, is:

\begin{lstlisting}[caption={Canonical training configuration (production model).},captionpos=b]
model:
  architecture: resnet50_softcon
  hub_name: SENTINEL2_ALL_SOFTCON
  pretrained_band_order: [B1,B2,B3,B4,B5,B6,B7,B8,B8A,B9,B10,B11,B12]
  num_classes: 4
  input_channels: 9
  freeze_backbone_epochs: 5
  class_names: [NotFarm, Poultry, Pigs, Cattle]
training:
  epochs: 50
  batch_size: 32
  learning_rate: 0.0001
  weight_decay: 0.01
  scheduler: cosine
  early_stopping_patience: 10
  checkpoint_metric: val_f1
  normalization: per_channel
  crop_center_px: 64
  balanced_class_sampling: false
  upsample_minority_regions: false
  class_weight: null
  seed: 42
  mixed_precision: true
inference:
  threshold: 0.5
  confidence_tiers: {high: 0.9, medium: 0.7, low: 0.4}
\end{lstlisting}

Optimisation uses AdamW \citep{loshchilov2019adamw} with cosine annealing
\citep{loshchilov2017sgdr}\unablated{E1.4}. The backbone is frozen for the first
five epochs while the fresh head converges; at unfreezing, the optimiser is
rebuilt with the learning rate scaled by $0.1$ and a fresh cosine schedule.
Model selection maximises validation macro-F1 rather than minimising validation
loss, on the grounds that loss is dominated by the majority class in a
50:1-imbalanced problem\unablated{E1.7}. Mixed precision is enabled for training
but deliberately \emph{disabled} at inference, because half-precision perturbs
probabilities at the $10^{-3}$ level and can flip \code{argmax} on near-ties,
breaking bit-comparability with published releases.

The loss is cross-entropy. Optional variants---class weighting, focal loss
\citep{lin2017focal}, and train-time logit adjustment
\citep{menon2021logit}---are implemented and evaluated in
Section~\ref{sec:imbalance}; the production model uses none of them.

Every run is evaluated after training by reloading the best checkpoint and
scoring each held-out slice separately, with per-country breakdowns written for
the \textsc{eval}, \textsc{generalization}, and \textsc{qual\_eval} slices. All
runs use a single seed\unablated{E0.3}; we have no variance estimate over seeds
anywhere in the experimental record, which is a material weakness we return to
in Section~\ref{sec:measurement}.

\subsection{Inference and deployment}
\label{sec:inference}

Scoring loads the patch metadata table, filters to the current imagery hash,
validates patch locations (below), reconstructs the exact training transform,
and runs the model. For the four-class model the prediction is \code{argmax} and
the full probability vector is persisted, which allows any downstream
binary farm/not-farm analysis to be recomputed as \(1 - P(\textsc{NotFarm})\)
without rescoring. Confidence tiers at $0.9/0.7/0.4$ support triage in the
review interface. Test-time augmentation over the eight dihedral transforms is
implemented but has never been enabled in a published run\unablated{E1.8}.

Full-world scoring is sharded by country across compute nodes using a greedy
longest-processing-time partition, with each shard writing to a distinct output
path; a merge step recombines them.

\paragraph{Patch-location validation.}
Before any training or scoring run, every patch is checked against the
\emph{current} coordinates of its candidate: the haversine distance between the
coordinates at which the patch was extracted and the candidate's present
centroid must be under \SI{250}{\metre}. Rows exceeding it are dropped, and if
more than 5\% of rows are stale the run aborts rather than proceeding. This
guard exists because of the incident described in Section~\ref{sec:geofix}, and
it is the single most important piece of defensive code in the pipeline.

\subsection{Geometric baseline}
\label{sec:if}

A parallel, purely geometric model provides a baseline that uses no imagery. An
Isolation Forest \citep{liu2008isolation,pedregosa2011sklearn} (500 estimators,
contamination $0.05$) is fitted on seven cluster-level morphometric features:
median building area, median length, number of unique orientations, pairwise
morphological distance, water area, mean built probability, and water proportion.
It is fitted \emph{only} on the \num{1200}-row Delmarva reference set, making it a
one-class template-similarity model: it scores how closely a cluster resembles a
clean United States broiler complex.

This baseline is informative because it isolates what geometry alone can do. On
\num{15489} labelled clusters, in binary farm-versus-not terms
(Table~\ref{tab:ifcnn}), it reaches ROC-AUC $0.862$ against the convolutional
model's $0.963$. The operational gap is larger than the AUC gap suggests: at
recall $0.95$, the Isolation Forest achieves precision $0.818$ while the
convolutional model achieves $0.959$, removing 86\% of non-farms against the
Isolation Forest's 30\%. A convex blend ($0.2$ geometric, $0.8$ convolutional)
improves average precision slightly to $0.9884$; using the Isolation Forest as a
\emph{sequential prefilter} before the convolutional model is strictly dominated
by simply thresholding the convolutional model. We therefore retain it as a
passthrough column and a diagnostic, not as a pipeline stage.

\begin{table}[t]
\centering
\small
\caption{Geometric baseline versus the convolutional model, binary
farm-versus-not, on \num{15489} labelled clusters. The always-predict-farm row
is the base rate implied by the candidate-generation stage and is the correct
reference point for all precision figures in this paper.}
\label{tab:ifcnn}
\begin{tabular}{lrrrr}
\toprule
Model & ROC-AUC & Avg.\ precision & Precision @ recall $0.95$ & Non-farms removed \\
\midrule
Always predict farm      & ---    & 0.767  & 0.767 & 0\% \\
Isolation Forest (geometry) & 0.862 & 0.952 & 0.818 & 30\% \\
CNN (imagery)            & 0.963  & 0.979  & 0.959 & 86\% \\
Blend ($0.2$/$0.8$)      & ---    & 0.988  & ---   & --- \\
\bottomrule
\end{tabular}
\end{table}

% =====================================================================
\section{Experiments and results}
\label{sec:results}
% =====================================================================

We report two experimental series. The first varies the \emph{model} (backbone,
context, imbalance treatment) on a three-class taxonomy
(\textsc{NotFarm}/\textsc{Poultry}/\textsc{OtherFarm}). The second holds the
architecture byte-identical and varies the \emph{labels}, on the four-class
taxonomy; the production model comes from this series. Unless noted, three-class
numbers are macro-F1 and four-class numbers are binary farm-versus-not ROC-AUC,
the latter because it is threshold-free and comparable across slices with very
different class balance.

\subsection{Backbone and spatial context}
\label{sec:backbone}

Table~\ref{tab:backbone} reports the backbone matrix. Each row differs from the
SSL4EO reference in exactly one respect, and the data and patch configuration
are byte-identical across rows.

\begin{table}[t]
\centering
\small
\caption{Backbone and context ablation, three-class macro-F1. \textsc{eval} is
in-domain held-out; \textsc{gen} is unseen countries. The ImageNet 9-channel row
is the control that separates pretraining from band count. Caveat: these are
single-seed point estimates on the contaminated patch store
(Section~\ref{sec:geofix}); see Table~\ref{tab:geofix} for corrected magnitudes
and the text for the confidence-interval analysis.}
\label{tab:backbone}
\begin{tabular}{llrrr}
\toprule
Run & Initialisation / change & \textsc{test} & \textsc{eval} & \textsc{gen} \\
\midrule
\code{v8\_three\_class}    & ImageNet, 4 channels          & 0.626 & 0.403 & 0.397 \\
\code{v9\_imagenet9ch}     & ImageNet, 9 channels          & 0.617 & 0.393 & 0.382 \\
\code{v8\_ssl4eo}          & SSL4EO-S12 MoCo, 9 channels   & 0.723 & 0.489 & 0.407 \\
\code{v9\_softcon}         & SoftCon, 9 channels           & 0.724 & \textbf{0.504} & 0.402 \\
\code{v9\_ctx128}          & SSL4EO, 128\,px context       & 0.729 & 0.484 & \textbf{0.440} \\
\code{v9\_softcon\_ctx128} & SoftCon $+$ 128\,px context   & \textbf{0.751} & 0.492 & 0.378 \\
\code{v8\_crt}             & SSL4EO $+$ cRT head           & 0.727 & 0.494 & 0.426 \\
\code{v8\_logitadj}        & SSL4EO $+$ logit-adjusted CE  & 0.572 & 0.470 & 0.330 \\
\bottomrule
\end{tabular}
\end{table}

\paragraph{Pretraining, not band count, drives the win.}
Sentinel-2--native pretraining improves \textsc{eval} macro-F1 from $0.403$ to
$0.489$--$0.504$. The obvious confound is that the ImageNet baseline used four
channels while the self-supervised runs used nine, so the gain might be extra
spectral information rather than better initialisation. The
\code{v9\_imagenet9ch} run resolves this: holding the nine-channel input fixed
and reverting only the initialisation gives $0.393$ \textsc{eval}---statistically
indistinguishable from the four-channel ImageNet run at $0.403$, and far below
every Sentinel-2--native run. The additional bands are close to worthless to an
ImageNet-initialised network and valuable to a Sentinel-2--pretrained one, which
is the expected behaviour if the pretrained filters encode multispectral
structure the ImageNet filters cannot represent. We regard this as the single
most robust experimental result in the project.

\paragraph{Everything else is a tie.}
A validity audit re-ran these comparisons after de-duplicating candidates (the
patch store can hold multiple rows per candidate, which inflated every split)
and attaching bootstrap confidence intervals \citep{efron1987bootstrap}.
De-duplication alone drops the SoftCon \textsc{eval} figure from $0.504$ to
$0.469$. After correction, the six leading Sentinel-2--native configurations
fall in the range $0.466$--$0.470$ with confidence-interval half-widths of
roughly $\pm0.05$: a statistical tie. The SoftCon-versus-SSL4EO difference of
$0.015$ is well inside noise. Only the Sentinel-2--native versus ImageNet
contrast, at roughly $+0.10$, exceeds twice the interval width. The production
choice of SoftCon over SSL4EO was therefore not justified by these measurements;
it rested on external benchmark evidence \citep{lacoste2023geobench}.

Re-testing the pair on the frozen benchmark confirms the production choice, but
by a far smaller margin than four-class macro-F1 suggests. SSL4EO is worse by
$0.0053$ ROC-AUC ($6.9\sigma$) and by $0.037$ macro-F1 on \textsc{qual\_eval},
and is indistinguishable from SoftCon on \textsc{test}, \textsc{eval}, and
\textsc{generalization} once the \textsc{Cattle} class is excluded. Raw
four-class macro-F1 instead reports SSL4EO as $0.143$ worse on \textsc{test},
which is an artefact: per class the two differ by $0.001$, $0.006$, and $0.010$
on \textsc{NotFarm}, \textsc{Poultry}, and \textsc{Pigs}, and by $0.541$ on
\textsc{Cattle}, whose F1 varies by up to $0.2$ between seeds of the same
configuration (Section~\ref{sec:measurement}).

\paragraph{Context does not compose.}
Increasing the crop from 64 to 128 pixels (\SI{640}{\metre} to
\SI{1.28}{\kilo\metre}) improves out-of-domain macro-F1 for SSL4EO
($0.407 \rightarrow 0.440$). A larger blind slice was previously reported to show
a $+0.047$ macro-F1 advantage for the wider-context model, but that result does
not survive re-testing: on the \num{1072} blind rows we can pair exactly across
the two models, the difference is $+0.002$ pooled (95\% CI $[-0.044, +0.048]$,
$p = 0.93$) and $+0.017$ when averaged over countries to neutralise composition
skew.

A controlled ablation settles the question against larger context. Training the
production configuration at 48, 64, and 128-pixel crops and evaluating on the
frozen benchmark, 128 pixels is \emph{worse} than the production 64
($-0.0044$ ROC-AUC, $5.7\sigma$ against the measured seed noise), as is 48
($-0.0026$, $3.4\sigma$). The per-slice pattern is the informative part: 128
pixels gains on \textsc{test} ($+0.015$ macro-F1) while losing heavily on
\textsc{generalization} ($-0.085$, $8.6\sigma$). Additional context appears to
help the model memorise the regions it trained on and to hurt transfer to unseen
ones, which is the opposite of what a screening deployment needs. The production
64-pixel crop is retained. Combining SoftCon with 128-pixel
context does \emph{not} stack: \textsc{gen} falls to $0.378$, below either lever
alone. We do not have a mechanistic explanation, and note that the production
model retains the 64-pixel crop largely for historical reasons. Given that
larger context wins on the two most deployment-like slices, this is among the
highest-value re-tests in the experiment plan\unablated{E1.2}.

\subsection{Class-imbalance interventions}
\label{sec:imbalance}

We evaluated five distinct mechanisms. None is used in the production model, and
the reasoning for each rejection differs.

\emph{Class-weighted cross-entropy} was tried with weights $[1.0, 0.7, 2.0]$ and
had no measurable effect---though this run coincided with the split bug of
Section~\ref{sec:validity} and is uninterpretable.

\emph{Classifier retraining} (cRT) \citep{kang2020decoupling}, freezing the
backbone and retraining the head with balanced sampling, improved out-of-domain
macro-F1 by $+0.019$ on the SSL4EO backbone but \emph{reduced} it by $-0.020$ on
SoftCon. A sign flip of this size on $n=273$ is inside sampling noise; the
honest reading is that cRT's effect here is unmeasured, not that it is
backbone-dependent.

\emph{Train-time logit adjustment} \citep{menon2021logit} reached
\textsc{eval} $0.470$ with minority-class F1 comparable to the unadjusted
baseline, but at a $0.15$ cost in \textsc{test} macro-F1. A post-hoc
$\tau$ sweep over the trained models found the optimum at $\tau = 0$ for the
best models, i.e.\ no adjustment---indicating the limiting factor is the
representation, not the decision rule.

\emph{Adaptive batch normalisation} \citep{li2018adabn}, recomputing
normalisation statistics on target-domain data, \emph{hurt}: out-of-domain
macro-F1 fell from $0.407$ to $0.383$. This is a useful negative: it suggests
the out-of-domain gap is not radiometric domain shift, which AdaBN addresses,
but morphological and label-distributional difference, which it does not.

\emph{Class-balanced sampling} is the intervention we studied most carefully,
because it appeared to be a win under contaminated data and then reappeared as a
trade under clean data. Measured twice on independent experiments, it improves
\textsc{eval} by $+0.027$ and $+0.046$ while reducing \textsc{generalization} by
$-0.016$ and $-0.054$. The mechanism is legible: \textsc{eval} is drawn from the
same training countries whose composition the sampler reweights, whereas
\textsc{generalization} contains no training rows at all, so upweighting the
minority classes optimises the slice that shares the training distribution at
the expense of the slice that does not. The clearest diagnostic is the
\textsc{Cattle} class: balanced sampling raises \textsc{test} \textsc{Cattle} F1
from $0.556$ to $0.636$ while dropping \textsc{qual\_eval} \textsc{Cattle} F1
from $0.250$ to $0.000$---memorisation of 153 training examples rather than a
learned category. Farm recall on inference countries falls from 89.3\% to 78.3\%.
Since deployment targets countries with no training data, we keep balancing off.

We tested whether this trade is instead an artefact of the spatial leakage
documented in Section~\ref{sec:measurement}---if the in-domain gain came from
clusters neighbouring training data, it would shrink once those rows were
removed. It does not: restricting \textsc{eval} to clusters at least
\SI{1.28}{\kilo\metre} from any training cluster \emph{widens} the gain from
$+0.046$ to $+0.084$ macro-F1, while on the leaked rows the balanced model is
$0.092$ \emph{worse}. The out-of-domain loss is likewise undiminished by
filtering ($-0.054$ macro-F1; $-0.019$ ROC-AUC, $p = 0.038$). The trade is a
genuine property of the intervention.

Country-level rebalancing is disabled permanently for a structural reason: 26
training countries have two or fewer rows, and inverse-frequency country weights
would draw a single-row country roughly 100 times per epoch.

\subsection{Label acquisition}
\label{sec:labelrounds}

Holding the architecture and recipe byte-identical, we varied only the label
set. Table~\ref{tab:labelrounds} reports binary farm-versus-not ROC-AUC per
slice.

\begin{table}[t]
\centering
\small
\caption{Label-round series, binary farm-versus-not ROC-AUC. Architecture and
training recipe are byte-identical across columns; only the label set differs.
Bracketed figures are 95\% bootstrap confidence intervals. \code{v9} is the
production model; \code{v10} is the rejected regularisation bundle
(Section~\ref{sec:v10}).}
\label{tab:labelrounds}
\begin{tabular}{lrrrrrr}
\toprule
Slice & $n$ & v6 & v7 & v8 & v9 & v10 \\
\midrule
\textsc{test}              & \num{2094}  & 0.991 & 0.993 & 0.993 & \textbf{0.994} [.989,.997] & 0.993 \\
\textsc{eval}              & 523         & 0.903 & 0.903 & \textbf{0.920} [.888,.949] & 0.911 & 0.903 \\
\textsc{generalization}    & 426         & 0.824 [.747,.894] & 0.843 & 0.909 & \textbf{0.915} [.885,.942] & 0.903 \\
\textsc{qual\_eval\_common} & \num{11437} & 0.966 & 0.954 & 0.981 & \textbf{0.983} [.980,.986] & 0.979 \\
ALB\,+\,COD (hardest OOD)   & 179         & 0.777 & 0.772 & 0.912 & \textbf{0.916} & --- \\
\bottomrule
\end{tabular}
\end{table}

\paragraph{Round~2: negatives alone shift the prior without improving the model.}
Round~2 added \num{962} adjudicated \textsc{NotFarm} corrections across 106
countries and nothing else. Superficially it looked like a large success: the
false-positive rate on the hardest out-of-domain pair fell from 66.2\% to 5.0\%.
But recall collapsed in step, from 92.5\% to 42.5\%. Three diagnostics together
establish that no ranking improvement occurred. ROC-AUC was flat or slightly
worse ($0.777 \rightarrow 0.772$ on ALB\,+\,COD; $0.979 \rightarrow 0.910$ on
BGD\,+\,NGA). At a \emph{matched} false-positive rate, recall was identical
(95.0\% for both models). And mean predicted farm probability fell for true
negatives ($0.091 \rightarrow 0.021$) \emph{and} for true positives
($0.846 \rightarrow 0.568$)---a uniform downward shift, not sharpened
separation. The model learned ``unfamiliar country implies not a farm.''
One-sided labels teach a prior, not a feature. This is the most transferable
methodological lesson in our record: any label campaign that adds only one class
should be evaluated with threshold-free and matched-operating-point metrics,
because threshold-dependent metrics will show a spurious improvement.

\paragraph{Round~3: balanced promotion is what actually helped.}
Adding \num{1740} farm \emph{positives} from the adjudicated pool, capped per
country, improved ranking on every slice: ALB\,+\,COD $0.772 \rightarrow 0.912$,
BGD\,+\,NGA $0.910 \rightarrow 0.936$, \textsc{qual\_eval} $0.952 \rightarrow
0.981$. Against the round-1 baseline, the \textsc{qual\_eval} false-positive rate
fell from 4.9\% to 3.6\% \emph{with} recall rising from 86.8\% to 88.9\%---a
genuine Pareto improvement, in contrast to round~2's trade. Per-class,
\textsc{Pigs} F1 gained $+0.124$.

\paragraph{The per-country cap is not a sensitive knob.}
Raising the cap from 50 to 70 promoted a further \num{1963} rows and changed
essentially nothing: ALB\,+\,COD $0.912 \rightarrow 0.916$, BGD\,+\,NGA
$0.936 \rightarrow 0.937$, \textsc{qual\_eval} unchanged at $0.983$,
\textsc{test} $0.993 \rightarrow 0.994$. Re-testing this on a purpose-built blind
benchmark of \num{11365} rows across 131 countries (Section~\ref{sec:measurement})
converts that observation into a properly powered null result: the paired
difference between the two label rounds is $-0.0003$ ROC-AUC with a 95\%
confidence interval of $[-0.0019, +0.0012]$ and $p = 0.715$, which excludes any
effect larger than about two thousandths of a point. On the same benchmark the
two earlier rounds separate decisively ($-0.014$ and $-0.031$, both $p < 0.001$),
confirming the benchmark can detect real differences of this kind. This lever is
exhausted at the current scale: further promotion from the same adjudicated pool
is unlikely to help, and the next data investment should change the label
\emph{distribution} (new regions, new species) rather than its volume.

\subsection{Regularisation bundle}
\label{sec:v10}

A final experiment combined label smoothing at $0.05$
\citep{szegedy2016rethinking}, cutout with two 16-pixel holes at probability
$0.5$ \citep{devries2017cutout}, and a longer schedule (70 epochs, patience 15).
It was run as a bundle, not as single levers---a design choice we now regard as
a mistake, since the result cannot be attributed.

The bundle was rejected. Label smoothing \emph{worsened} measured calibration,
which is the opposite of its usual justification: expected calibration error
\citep{guo2017calibration} rose from $0.027$ to $0.052$ on
\textsc{qual\_eval\_common} and from $0.006$ to $0.020$ on \textsc{test}. Cutout
did not achieve its stated purpose of improving the data-poor \textsc{Cattle}
class---one-versus-rest AUC for \textsc{Cattle} on \textsc{test} fell from
$0.936$ to $0.856$, and the out-of-domain \textsc{Cattle} slice remained at zero
correct out of four. Cutout is also applied to the full 128-pixel patch before
the 64-pixel centre crop, so a hole lands inside the visible crop only about a
quarter of the time per hole, making the effective augmentation rate roughly 22\%
of samples rather than the nominal 50\%. The one gain, $+0.057$ macro-F1 on
\textsc{eval}, is on the slice that deployment does not target, and an ensemble
of the bundle with the production model performed identically to the production
model alone, indicating no added diversity.

The broader conclusion we drew is that training-recipe tweaks on the current
label set are exhausted, and that further effort belongs in measurement and data.

\subsection{Controlled ablation of the inherited defaults}
\label{sec:ablations}

The design choices in Section~\ref{sec:methods} were, for most of this project's
history, inherited defaults rather than measured decisions. We tested them
directly: twenty runs, each the production configuration with exactly one change,
evaluated on the frozen benchmark of Section~\ref{sec:measurement} and judged
against the seed noise measured on the same benchmark ($\sigma = 0.0008$).
Table~\ref{tab:ablations} reports the result.

\begin{table}[t]
\centering
\small
\caption{Single-lever ablations, farm-versus-not ROC-AUC on the frozen blind
benchmark ($n = \num{11365}$), as a difference from the production
configuration. $\Delta/\sigma$ uses the seed standard deviation measured on the
same benchmark. Exactly one change improves the model.}
\label{tab:ablations}
\begin{tabular}{llrrl}
\toprule
Experiment & Change & $\Delta$AUC & $\Delta/\sigma$ & Verdict \\
\midrule
E1.4 & \textbf{no backbone freeze phase} & $\mathbf{+0.0045}$ & $+5.9$ & \textbf{better} \\
E1.5 & drop photometric augmentations & $+0.0010$ & $+1.3$ & tie \\
E1.1 & recompute indices after jitter & $+0.0009$ & $+1.2$ & tie \\
E1.4 & learning rate $3\times10^{-4}$ & $+0.0007$ & $+0.9$ & tie \\
E1.5 & cutout only & $+0.0004$ & $+0.5$ & tie \\
E1.7 & select on validation loss & $+0.0000$ & $+0.0$ & tie \\
E1.9 & three-class taxonomy & $+0.0000$ & $+0.0$ & tie \\
E1.4 & learning rate $3\times10^{-5}$ & $-0.0004$ & $-0.5$ & tie \\
E1.1 & \textbf{six bands, no indices} & $-0.0010$ & $-1.3$ & \textbf{tie} \\
E1.5 & geometric augmentations only & $-0.0022$ & $-2.8$ & worse \\
E1.2 & context 48\,px & $-0.0026$ & $-3.4$ & worse \\
E1.2 & context 128\,px & $-0.0044$ & $-5.7$ & worse \\
E1.6 & SSL4EO backbone & $-0.0053$ & $-6.9$ & worse \\
E1.1 & RGB + NIR & $-0.0102$ & $-13.4$ & worse \\
E1.1 & RGB + NDWI & $-0.0115$ & $-15.2$ & worse \\
E1.1 & RGB only & $-0.0170$ & $-22.3$ & worse \\
\bottomrule
\end{tabular}
\end{table}

\paragraph{The staged freeze hurts---but not for the reason it appears to.}
Disabling the five-epoch freeze phase improves the frozen benchmark by $+0.0045$
($5.9\sigma$, CI $[+0.0025, +0.0065]$), \textsc{test} macro-F1 by $+0.017$ and
\textsc{qual\_eval} by $+0.020$, while shortening training. This is the only
change of the twenty that improved the model, and we recommend adopting it.

The \emph{mechanism}, however, is not the one the lever's name implies. Inspection
of the training loop shows that the unfreeze transition does two things at once:
it unfreezes the backbone \emph{and} rebuilds the optimiser at
$0.1\times$ the learning rate, with a fresh cosine schedule. Setting
\texttt{freeze\_backbone\_epochs} to zero skips that branch entirely, so the
``no freeze'' arm also trains its backbone at $10\times$ the learning rate
of the baseline ($10^{-4}$ against $10^{-5}$) for every remaining epoch. We
confirmed this in the logged learning rate, which drops from
$9.755\times10^{-5}$ at epoch~5 to $9.990\times10^{-6}$ at epoch~6 in every
freeze-enabled run. The ablation is therefore a two-factor change reported as
one, and the second factor is very likely the dominant term.

This also qualifies the learning-rate result immediately below. The sweep varied
the \emph{base} rate while leaving the $0.1\times$ coupling intact, so the
configuration that a staged warm-up would most plausibly want---warm-up followed
by fine-tuning at the full rate---was never expressible, let alone tested. We
report the recipe-level finding as it stands and add an arm that decomposes it
in the round-four campaign of Section~\ref{sec:round4}: warm-up retained,
unfreeze at full learning rate. Until that arm reports, the correct statement is
that \emph{the freeze schedule's coupled learning-rate cut hurts}, not that
freezing hurts. The base learning rate is otherwise at a sensible optimum:
$3\times10^{-5}$ is clearly worse and $3\times10^{-4}$ trades \textsc{test}
against \textsc{qual\_eval}.

\paragraph{The spectral indices contribute nothing; the SWIR bands do.} Six raw
bands without NDVI, NDBI, or NDWI are indistinguishable from the full
nine-channel input on every slice ($-0.0010$, $1.3\sigma$), whereas reducing
further to RGB\,+\,NIR costs $13\sigma$. The information the model uses is in the
bands---including the two SWIR channels---and the three derived indices, a third
of the input, are dead weight. This is the expected result for quantities that
are deterministic functions of channels the network already sees, and it also
retires the index/jitter inconsistency noted in Section~\ref{sec:methods} as a
concern. Recomputing the indices after augmentation is likewise a no-op.

\paragraph{Augmentation earns its place, but only out of domain.} Reducing the
recipe to flips and rotations costs both the frozen benchmark ($-0.0022$) and
\textsc{generalization} ($-0.036$ macro-F1, $3.6\sigma$); dropping only the four
photometric terms is a tie on the frozen benchmark but $-0.021$ on
\textsc{generalization}. The augmentations buy transfer, not in-domain accuracy.
Cutout in isolation is neutral everywhere, which retrospectively clears it of
responsibility for the failure of the regularisation bundle in
Section~\ref{sec:v10}.

\paragraph{Two settings do not matter.} The checkpoint-selection metric and the
class taxonomy are both exactly $0.0000$ on the frozen benchmark. The taxonomy
result is the more consequential: dropping to three classes costs nothing
measurable, so the \textsc{Cattle} class can be merged away at no cost---and
doing so would remove the dominant source of metric noise identified in
Section~\ref{sec:measurement}.

\subsection{Calibration and thresholds}
\label{sec:thresholds}

Miscalibration in this system is domain-conditional, which has a direct
operational consequence. Re-tuning a \emph{global} decision threshold is a
no-op: the validation-optimal farm threshold is $0.481$ against a default of
$0.5$, worth at most $2.5$ percentage points of out-of-domain recall. But
applying a \emph{lower threshold specifically out of domain} recovers
substantial recall at little false-positive cost, as the model is
systematically under-confident on unseen countries
(Table~\ref{tab:thresholds}). Pre-registering the rule ``adopt the lowest
threshold buying at least five points of recall for at most five points of
false-positive rate'' selects $0.4$ ($+0.055$ recall for $+0.018$ false-positive
rate); $0.3$ narrowly fails it at $+0.053$ false-positive rate. We therefore
recommend deployment with a domain-conditional threshold: the default in training
countries, and $0.4$ for screening in countries with no training data.

We also tested \emph{per-country} thresholds, fitted leave-one-country-out so
that no country's threshold sees its own rows, and they are clearly worse than a
single out-of-domain constant: pooled, they buy $+0.114$ recall at a cost of
$+0.170$ false-positive rate, against $+0.055$ for $+0.018$ from the fixed
threshold. With only four out-of-domain countries each fit rests on three others
and is unstable---one country's fitted threshold collapses to $0.05$ and triples
its false-positive rate. Per-country calibration is not viable until many more
labelled out-of-domain countries exist.

\begin{table}[t]
\centering
\small
\caption{Threshold sweep for the production model on the out-of-domain slice.
Lowering the threshold from $0.5$ to $0.3$ buys 9 points of recall for 5 points
of false-positive rate---a favourable trade when the output feeds human review.}
\label{tab:thresholds}
\begin{tabular}{lrrr}
\toprule
Threshold & Recall & False-positive rate & Precision \\
\midrule
0.2 & 0.894 & 0.211 & 0.864 \\
0.3 & 0.855 & 0.170 & 0.883 \\
0.4 & 0.820 & 0.135 & 0.901 \\
0.5 (default) & 0.765 & 0.117 & 0.907 \\
0.6 & 0.718 & 0.094 & 0.920 \\
\bottomrule
\end{tabular}
\end{table}

\subsection{Per-class performance}
\label{sec:perclass}

Table~\ref{tab:perclass} gives per-class results for the production model. Two
structural limits are visible. \textsc{Pigs} is systematically confused with
\textsc{Poultry}---the dominant error mode on \textsc{eval}, unmoved by every
intervention we tried---which is unsurprising at \SI{10}{\metre}, where the two
facility types differ mainly in barn dimensions and lagoon configuration.
\textsc{Cattle} is not measurable: 153 training rows, 25 in \textsc{test}, seven
in \textsc{eval}, four in \textsc{generalization}, six in \textsc{qual\_eval}.
Its out-of-domain one-versus-rest AUC of $0.485$ is indistinguishable from
chance, and any \textsc{Cattle} number in this paper should be read as a
placeholder awaiting labels rather than as a measurement.

\begin{table}[t]
\centering
\small
\caption{Per-class one-versus-rest AUC and F1 for the production model, with
support. \textsc{Pigs} is absent from the out-of-domain slice; \textsc{Cattle}
support is too small everywhere for its numbers to carry information.}
\label{tab:perclass}
\begin{tabular}{lrrrrrrrr}
\toprule
& \multicolumn{2}{c}{\textsc{test}} & \multicolumn{2}{c}{\textsc{eval}}
& \multicolumn{2}{c}{\textsc{gen}} & \multicolumn{2}{c}{\textsc{qual\_eval}} \\
\cmidrule(lr){2-3}\cmidrule(lr){4-5}\cmidrule(lr){6-7}\cmidrule(lr){8-9}
Class & AUC & F1 & AUC & F1 & AUC & F1 & AUC & F1 \\
\midrule
NotFarm & 0.994 & 0.923 & 0.911 & 0.752 & 0.915 & 0.789 & 0.983 & 0.976 \\
Poultry & 0.975 & 0.939 & 0.855 & 0.836 & 0.906 & 0.787 & 0.975 & 0.817 \\
Pigs    & 0.965 & 0.795 & 0.775 & 0.421 & \multicolumn{2}{c}{---} & 0.942 & 0.472 \\
Cattle  & 0.936 & 0.556 & 0.837 & 0.400 & 0.485 & 0.000 & 0.908 & 0.286 \\
\midrule
Support & \multicolumn{2}{c}{\num{2094}} & \multicolumn{2}{c}{523}
        & \multicolumn{2}{c}{426} & \multicolumn{2}{c}{\num{11437}} \\
\bottomrule
\end{tabular}
\end{table}

\subsection{Global deployment}
\label{sec:deployment}

The production model scores \num{157099} of \num{157102} candidates across 167
countries (Table~\ref{tab:tiers}), of which approximately 68\% are predicted to
be farms of some type---consistent with the 76.7\% farm base rate of the
candidate pool after removing predicted negatives.

Per-country out-of-domain AUC spans a wide range and is the most honest summary
of where the system can be trusted: South Africa $1.000$, Argentina $0.996$,
Portugal $0.996$, France $0.995$, Germany $0.981$, Australia $0.982$, Poland
$0.980$, Romania $0.979$, Hungary $0.973$, Italy $0.969$, Czechia $0.968$,
Bangladesh $0.947$, Albania $0.942$, Nigeria $0.917$, DR Congo $0.814$.
Performance degrades with distance from the training distribution in both the
geographic and the construction-practice sense: the weakest results are in
regions where facility architecture, roofing materials, and surrounding land
cover differ most from the North American and European examples that dominate
training. Separate analyses identified near-total failures in specific regions
(Russia, Ukraine, Malaysia, India, Turkey) where farm-class F1 falls below
$0.3$; these countries should be treated as unsupported by the current release.

Label-round changes also move deployment-scale outputs substantially. Over the
\num{154905} commonly scored points, predicted farms went from \num{113137}
(73.0\%) under round~1 to \num{87265} (56.3\%) under round~2 and \num{107013}
(69.1\%) under round~3---a reminder that the round~2 prior shift was not a
subtle statistical artefact but a change that would have removed roughly a
quarter of the world's predicted facilities had it shipped.

% =====================================================================
\section{Data integrity and measurement validity}
\label{sec:validity}
% =====================================================================

This section reports failures. We include it because two of them invalidated
published conclusions, and because the mechanisms are general enough to recur in
any pipeline that joins imagery to labels through an identifier.

\subsection{Patch-store key misalignment}
\label{sec:geofix}

\paragraph{Discovery.} A cluster in Mexico was scored as a farm with probability
$0.963$. Manual inspection confirmed a real egg-production facility at the
coordinates the \emph{patch} was extracted from---but the current label file
placed that cluster identifier several hundred kilometres away, on a
\textsc{NotFarm}.

\paragraph{Mechanism.} Patches were cached and joined by cluster identifier and
imagery hash, with no spatial check, and extraction skipped any identifier
already present in the store. An upstream re-run of the candidate-merge process
renumbered 93.5\% of cluster identifiers---geometry and labels were unchanged;
only the identifier-to-cluster mapping churned. From that moment, every join
silently paired a candidate with another candidate's imagery, and the skip-cache
guaranteed the wrong patch was never re-extracted.

\paragraph{Blast radius.} Measuring displacement between each patch's extraction
coordinates and its candidate's current coordinates showed \num{34.3}--34.6\% of
labelled-country rows displaced by more than \SI{250}{\metre}, with median
displacement in the hundreds of kilometres, against a $0.04\%$ control rate in
unaffected regions. Contamination was spread across \emph{all} splits, including
\textsc{eval} and \textsc{generalization}, at country-level rates of 21\%
(Thailand, Bangladesh) to 48\% (Chile). In one binary run, 949 of \num{2495}
rows were mislocated, and of the 384 for which the original location's label was
recoverable, 84 (22\%) were outright image/label contradictions. Every run
between the incident and the fix trained and evaluated on approximately
one-third wrong-image pairs; all metrics from that period are \emph{floors},
not estimates.

\paragraph{Correction and consequences.} We re-keyed the store by location,
re-extracted roughly \num{24700} stale patches, and reduced residual staleness to
$0.31\%$. Table~\ref{tab:geofix} shows the effect of re-running affected
configurations on clean data. The corrections are large---larger than any
modelling intervention in this paper---and they resolved a standing puzzle in
which an early model appeared to outperform its successors: the apparent
regression coincided exactly with the onset of contamination.

\begin{table}[t]
\centering
\small
\caption{Effect of correcting the patch-store misalignment, three-class
macro-F1. The correction is worth more than any modelling intervention we
tested, and it reversed two experimental verdicts.}
\label{tab:geofix}
\begin{tabular}{lrrrrrr}
\toprule
& \multicolumn{3}{c}{Contaminated} & \multicolumn{3}{c}{Corrected} \\
\cmidrule(lr){2-4}\cmidrule(lr){5-7}
Run & \textsc{test} & \textsc{eval} & \textsc{gen} & \textsc{test} & \textsc{eval} & \textsc{gen} \\
\midrule
\code{v9\_softcon} & 0.712 & 0.462 & 0.400 & 0.890 & 0.631 & 0.486 \\
\code{v9\_ctx128}  & 0.710 & 0.469 & 0.415 & 0.902 & 0.641 & 0.443 \\
\code{v3\_three\_class} & 0.526 & 0.399 & 0.298 & 0.716 & 0.523 & 0.397 \\
\bottomrule
\end{tabular}
\end{table}

Two verdicts reversed under clean data. Class-balanced sampling had been recorded
as ``the boundary moves but quality does not''; on corrected data it produced a
genuine $+0.027$ \textsc{eval} gain with \textsc{Pigs}/\textsc{Poultry} confusion
falling from 52\% to 38\%. Larger context had been recorded as non-additive with
SoftCon; on corrected data it gave a partial $+0.017$. Neither reversal would
have been detectable without re-running.

\paragraph{Invariants now enforced.} Three code-level rules resulted. First,
patch locations are validated against current candidate coordinates before every
run, with a hard abort above 5\% staleness (Section~\ref{sec:inference}).
Second, the candidate-merge step asserts identifier stability---at least 99.9\%
overlap and centroid drift under $10^{-6}$ degrees---and refuses to join
otherwise. Third, upstream cluster identifiers are treated as \emph{ephemeral}:
they may be used within a single data pull but never to join across pulls.
Location, not identity, is the durable key.

\subsection{Label-source shift}
\label{sec:labelshift}

The gap between \textsc{test} and \textsc{eval} performance is not primarily a
generalisation gap; it is a change in what the labels mean.

For the minority farm class, training labels are 98\% registry-derived and 56\%
United States, while evaluation labels are 80\% visually adjudicated and 3\%
United States. The two are different measurement instruments applied to
different populations. The same holds for the negative class: the large majority
of training \textsc{NotFarm} rows are pipeline-generated negatives with no
recorded label source, whereas nearly all evaluation \textsc{NotFarm} rows are
visually confirmed hard negatives---structures a human considered plausibly
farm-like before rejecting them. \textsc{NotFarm} F1 falls from $0.73$ on
\textsc{test} to $0.50$ on \textsc{eval} largely for this reason.

Holding the split fixed and varying only the labelling instrument makes the
effect directly visible. Within \textsc{eval}, the production model scores
ROC-AUC $1.000$ on the 73 registry-labelled rows and $0.909$ on the 449
visually-labelled ones; the headline \textsc{eval} figure of $0.911$ is therefore
essentially the visual-label figure, while \textsc{test}, at 84\% registry, is
essentially the registry figure. The same ordering holds in every split
(\textsc{test} $-0.043$ ROC-AUC, \textsc{qual\_eval} $-0.019$), and it is
near-zero within \textsc{train} itself ($-0.002$): the model fits both label
sources equally well during training and generalises worse only on visual ones.
That pattern points to visually-adjudicated rows being intrinsically harder
cases---hard negatives and ambiguous sites---rather than to registry labels being
noisy, which argues for acquiring more visual labels for training rather than
down-weighting registry ones.

The consequence is that \textsc{test}-to-\textsc{eval} degradation should not be
read as overfitting, and interventions aimed at overfitting will not close it.
The dominant evaluation error mode---minority farm types predicted as
\textsc{Poultry}---persists at similar rates whether the evaluation row was
registry-labelled (18\% correct) or visually labelled (24\% correct), which
indicates a genuine representational limit rather than a pure annotation
artefact\unablated{E2.6}.

\subsection{Measurement limits}
\label{sec:measurement}

We close with the constraint that shaped this project's second half.

\paragraph{The benchmark is smaller than the effects.} At $n=523$, the 95\%
confidence interval on \textsc{eval} macro-F1 is approximately $\pm0.05$.
Detecting a difference $\Delta$ requires roughly $n \approx 524 \cdot
(0.05/\Delta)^2$: about \num{1455} rows for $\Delta = 0.03$, \num{3275} for
$\Delta = 0.02$, and \num{13100} for $\Delta = 0.01$. Nearly every intervention
in Section~\ref{sec:results} falls in the $0.01$--$0.05$ range. Most of our
experiments were, strictly, underpowered.

\paragraph{Splits are not spatially blocked.} Measuring the distance from every
held-out cluster to the nearest of the \num{12062} training clusters shows
substantial exposure: 16.6\% of \textsc{eval} clusters, 38.0\% of \textsc{test},
and 32.5\% of \textsc{val} lie within one patch width
(\SI{1.28}{\kilo\metre}) of a training cluster, so their imagery partially
overlaps training imagery. The spatial cross-validation literature
\citep{roberts2017spatialcv,ploton2020spatial} is unambiguous that this inflates
apparent performance, and it does here: splitting each slice at the
\SI{1.28}{\kilo\metre} boundary, the production model scores ROC-AUC $0.961$ on
near \textsc{eval} rows against $0.905$ on far ones ($+0.056$), and $0.993$
against $0.964$ on \textsc{val} ($+0.028$). All in-domain figures in this paper
should be read as optimistic by roughly that margin. Two slices are exempt:
\textsc{generalization} has effectively no exposure (0.5\% within
\SI{1.28}{\kilo\metre}, median distance \SI{116}{\kilo\metre}) and
\textsc{qual\_eval} has 3.5\%, which is why we weight those two most heavily.
This measurement bounds evaluation contamination only; the model still trained on
the nearby clusters, so a spatially blocked \emph{retrain} is still
required\unablated{E0.2}.

\paragraph{Duplicates inflated every split.} The patch store can hold multiple
rows per candidate; de-duplication reduces \textsc{eval} macro-F1 by
approximately $0.03$. Figures in Table~\ref{tab:backbone} are pre-de-duplication
point estimates and are comparable to one another but not to post-correction
figures.

\paragraph{Seed variance is larger than most measured effects.} We trained the
production configuration five times, varying only the random seed. On the frozen
benchmark of Section~\ref{sec:measurement} the model is highly reproducible
(farm ROC-AUC $\sigma = 0.0008$). On the 523-row evaluation slice it is not:
four-class macro-F1 ranges from $0.602$ to $0.683$, giving $\sigma = 0.033$ and a
decision band of $\pm0.095$ for a comparison between two single runs.

That band is wider than nearly every intervention this project has ranked on that
metric. Class-balanced sampling ($+0.046$), the regularisation bundle
($+0.057$), and SoftCon over SSL4EO ($+0.015$) are all inside it. Those
comparisons were not measurable on the instrument used to measure them, and we
report them here as ties rather than as effects. Two conclusions survive on other
evidence: the balanced-sampling \emph{out-of-domain} loss is $-0.054$ where
$\sigma = 0.0075$, and the null result for the third label round is confirmed
independently on the frozen benchmark ($p = 0.715$).

\paragraph{The instability is a metric artefact, not model behaviour.}
Decomposing the seed spread by class identifies its source. On
\textsc{qual\_eval}, \textsc{NotFarm} F1 (9,676 rows) varies by $0.003$ across
seeds while \textsc{Cattle} F1 (six rows) varies by $0.150$. Because macro-F1
weights all classes equally, that single near-empty class contributes roughly
80\% of the variance in the headline number. The control is \textsc{test}, where
\textsc{Cattle} has 25 rows instead of six and the macro-F1 spread falls to
$0.014$. Excluding \textsc{Cattle} halves the evaluation-slice $\sigma$ to
$0.016$ and reduces the \textsc{qual\_eval} $\sigma$ threefold to $0.0056$.

Reporting four-class macro-F1 on slices where the rarest class has single-digit
support is therefore close to reporting a coin flip, and we recommend against it.
This is also a cautionary result for reading our own tables: raw macro-F1 made an
SSL4EO run appear $0.143$ worse than SoftCon on \textsc{test}---an apparent
twenty-five-sigma collapse---when the three measurable classes differed by
$0.006$ and the entire gap was \textsc{Cattle} falling from $0.541$ to zero in a
single run.

\paragraph{A blind benchmark closes most of this gap.} The measurement problem is
largely solvable with data already in hand. We froze a benchmark of \num{11365}
labelled clusters spanning 131 countries, drawn from the adjudicated
\textsc{qual\_eval} pool and further filtered to clusters at least
\SI{1.28}{\kilo\metre} from any training cluster, so that it is blind in both
senses: no model trained on these rows, and none is a spatial near-duplicate of a
training row. Its 95\% bootstrap confidence half-width is $0.0030$ on farm
ROC-AUC and $0.0061$ on macro-F1---roughly sixteen times tighter than the
$\pm0.05$ of the 523-row evaluation slice, and enough to make differences of
$0.01$ comfortably detectable. Two results in this paper were re-derived on it
(Sections~\ref{sec:backbone} and~\ref{sec:labelrounds}), and it should be the
default slice for future comparisons. Its own limitation is composition: it is
dominated by \textsc{NotFarm} rows and contains only six \textsc{Cattle}
examples, so it resolves farm-versus-not questions well and minority-type
questions not at all.

% =====================================================================
\section{The round-four campaign: recipe levers under seed replication}
\label{sec:round4}

After the analyses above were completed, a fourth labeling round restructured
the training data: every row not reserved for the generalization hold-out was
assigned to train or validation, absorbing the qualitative-evaluation slice
entirely. Training data grew from 12{,}062 to 21{,}478 rows (+73\%), validation
from 2{,}666 to 5{,}022, while the test (2{,}094), eval (523) and generalization
(617 usable; 662 labeled) hold-outs kept their membership. The frozen benchmark
of Section~\ref{sec:measurement} was thereby retired---its rows are now training
data---so the campaign relies on the three surviving hold-outs, for which we
verified 0.0\% overlap with the train and validation sets of every archived
model (v6--v9), and that each archived run trained on its recorded explicit
splits. The campaign asks four pre-registered questions: did the new labels
move performance; do the two recipe changes recommended by
Section~\ref{sec:ablations} (no freeze phase; six bands) replicate and compose;
and how does an ImageNet-initialized DenseNet-121 compare.

\paragraph{Configuration.} All arms share the production recipe of
Table~\ref{tab:r4recipe} except for the single lever each one changes
(Table~\ref{tab:r4arms}). Every arm was trained three times, varying only the
random seed $\in\{42,43,44\}$, which controls the initialization of the fresh
four-class head, the per-epoch shuffle order, and every augmentation draw;
data, splits and imagery are byte-identical across all 18 runs (verified at the
level of candidate-id sets).

\begin{table}[t]
\centering
\small
\caption{The shared training configuration (arm A is exactly this recipe, which
is byte-identical to the v9 production configuration).}
\label{tab:r4recipe}
\begin{tabular}{ll}
\toprule
Component & Setting \\
\midrule
Backbone & ResNet-50, SoftCon Sentinel-2 pretraining \citep{wang2024softcon} \\
First convolution & band-mapped from the 13-band pretrained stem \\
Head & fresh linear layer, 4 classes (NotFarm/Poultry/Pigs/Cattle) \\
Input & 9 channels (B2,B3,B4,B8,B11,B12 + NDVI,NDBI,NDWI), $64{\times}64$\,px crop \\
Normalization & per-channel, training-set statistics \\
Optimizer & AdamW \citep{loshchilov2019adamw}, lr $10^{-4}$, weight decay $0.01$ \\
Schedule & cosine \citep{loshchilov2017sgdr}, $T_{\max}=50$ epochs, batch size 32 \\
Freeze phase & backbone frozen 5 epochs; optimizer rebuilt at $0.1\times$ lr on unfreeze \\
Early stopping & patience 10 on validation macro-F1; best-val checkpoint shipped \\
Augmentation & horizontal/vertical flips ($p{=}0.5$), $90^\circ$ rotations, scale jitter, cutout \citep{devries2017cutout} \\
Class handling & no balancing, no class weights; mixed precision \\
\bottomrule
\end{tabular}
\end{table}

\begin{table}[t]
\centering
\small
\caption{Arms. B--D and F are single- or double-lever deltas against A;
arm F was added mid-campaign after code inspection revealed that the unfreeze
transition rebuilds the optimizer at $0.1\times$ the learning rate, so that
``no freeze'' had always changed two factors at once. F keeps the warm-up and
removes only the learning-rate cut---a configuration that was previously
unrepresentable. The amendment predates any arm-F result and made the
Holm correction stricter for the original contrasts.}
\label{tab:r4arms}
\begin{tabular}{lll}
\toprule
Arm & Delta vs.\ A & Question \\
\midrule
A & none & label-round effect (vs.\ archived v9) \\
B & \code{freeze\_backbone\_epochs: 0} & does the freeze-phase result replicate? \\
C & 6 bands (drop NDVI/NDBI/NDWI) & do the indices earn their place? \\
D & B $+$ C & do the levers compose? \\
E & D $+$ DenseNet-121 (ImageNet init, 7.0M vs.\ 23.5M params) & architecture family \\
F & \code{unfreeze\_lr\_scale: 1.0} (freeze kept) & freeze vs.\ learning rate, decomposed \\
\bottomrule
\end{tabular}
\end{table}

\paragraph{Protocol.} The primary estimand is recipe-level: the mean of
per-seed AUCs, never a selected seed and never the seed-ensemble. Every
contrast's standard error carries a mandatory seed term,
$\mathrm{SE}^2 = \mathrm{SE}^2_{\text{boot}} +
\sigma^2_{\text{seed}}(1/n_1{+}1/n_2)$, with $\sigma_{\text{seed}}$ pooled per
slice; the five confirmatory contrasts are Holm-corrected
\citep{holm1979simple} and additionally judged against a practical floor of
0.005 AUC. Seed replication is not optional in this regime: seed sensitivity
of fine-tuned pretrained models is well documented
\citep{dodge2020finetuning,picard2021seed}, and this campaign reproduced its
canonical failure mode live---with one collected seed, arm B ``beat'' A at
$+0.050$ AUC ($p=0.003$); with all three, the same contrast is $+0.023$ at
$p=0.29$.

\begin{table}[t]
\centering
\small
\caption{Farm-vs-not ROC-AUC by slice: three-seed mean $\pm$ sd per arm;
archived models scored on identical rows. Generalization $n{=}632$, eval
$n{=}572$, test $n{=}2{,}162$ (rows labeled \emph{Ambiguous} are excluded from
the binary target; farm-type-unknown labels count as positives).}
\label{tab:r4results}
\begin{tabular}{lccc}
\toprule
Model & Generalization & Eval & Test \\
\midrule
A (baseline) & $0.833 \pm 0.009$ & $0.915 \pm 0.004$ & $0.993 \pm 0.000$ \\
B (freeze0) & $0.857 \pm 0.036$ & $0.929 \pm 0.004$ & $0.995 \pm 0.001$ \\
C (6 bands) & $0.833 \pm 0.005$ & $0.917 \pm 0.004$ & $0.992 \pm 0.001$ \\
D (B$+$C) & $0.861 \pm 0.033$ & $0.933 \pm 0.006$ & $0.994 \pm 0.001$ \\
E (DenseNet) & $0.815 \pm 0.012$ & $0.934 \pm 0.010$ & $0.991 \pm 0.001$ \\
F (full-lr unfreeze) & $0.850 \pm 0.027$ & $0.929 \pm 0.007$ & $0.995 \pm 0.001$ \\
\midrule
archived v9 (round 3) & 0.874 & 0.906 & 0.993 \\
\bottomrule
\end{tabular}
\end{table}

\paragraph{Results.} No confirmatory contrast survives correction on the
primary (generalization) slice: $b{>}a$ $+0.024$ ($p_{\text{Holm}}=0.67$),
$d{>}a$ $+0.028$ ($p=0.60$), $f{>}a$ $+0.017$ ($p=0.80$), $b{>}f$ $+0.008$
($p=0.80$), $e{>}d$ $-0.046$ ($p=0.16$); on the saturated test slice the only
nominally significant contrast ($e{>}d$, $-0.003$) is below the practical
floor. Three findings organize the campaign.

\emph{(i) The label-round effect is a trade.} Arm A isolates the data change:
against archived v9 it gains $+0.009$ on in-domain eval, is unchanged on test
($-0.000$), and loses $0.041$ on generalization. The added rows are the
absorbed qualitative-evaluation slice, concentrated in the focal countries;
denser in-domain data specialized the model at the expense of out-of-domain
transfer.

\emph{(ii) The freeze-phase result of Section~\ref{sec:ablations} decomposes
into a learning-rate effect with a reliability cost.} Arm F attributes
roughly two thirds of the freeze0 delta to the backbone learning rate
($f{>}a$, $+0.017$) and little to the warm-up itself ($b{>}f$, $+0.008$).
The higher rate also quadruples seed variance: arms training the backbone at
$10^{-5}$ show $\sigma_{\text{seed}} = 0.005$--$0.009$ on generalization,
arms at $10^{-4}$ show $0.027$--$0.036$, with one seed in three collapsing by
$\sim$0.06 AUC in each of B, D and F.

\emph{(iii) The variance is predominantly cross-country calibration drift,
not ranking skill.} Decomposing pooled generalization AUC into a
size-weighted mean of within-country AUCs, the collapsing seeds have normal
within-country skill (e.g.\ arm D's weakest seed: 0.847 within-country
against a pooled 0.823); what moves is the alignment of per-country score
levels with per-country prevalence. The collapsing seed's per-country offsets
correlate at 0.67--0.79 across arms B, D and F---the seed fixes the shuffle
and augmentation stream, a shared treatment whose imprint the higher learning
rate amplifies. Within countries the six arms are nearly equivalent
(0.832--0.850), while archived v9 reaches 0.868: the regression against v9 is
genuine ranking loss, and pooled cross-country AUC is the less stable and less
deployment-relevant of the two views. Calibration is uniformly poor
out-of-domain (ECE 0.16--0.20 \citep{guo2017calibration}, vs.\ $\sim$0.01
in-domain), so published scores function as rankings rather than
probabilities.

Secondary results: dropping the spectral indices replicates as a clean null
($c$ vs.\ $a$: $-0.000$ on generalization) with the lowest seed variance of
any arm, confirming Section~\ref{sec:ablations} on fresh labels; the DenseNet
arm is uninterpretable as an architecture verdict (it differs from D in
architecture, pretraining and capacity simultaneously) but its
profile---best on in-domain eval (0.934), last pooled out-of-domain
(0.815), ordinary within-country (0.846)---is the signature of an in-domain
specialist. Per-country results show the same two countries (Morocco, India;
AUC 0.65--0.76) weakest under every recipe: the between-country spread
exceeds every between-arm difference.

\paragraph{Consequences.} The round-three model remains the production model:
no round-four arm matches its out-of-domain ranking, pooled or within-country,
and the pre-registered retention rule held. Of the previously recommended
changes, the six-band simplification is adopted (equal accuracy, best
reproducibility, three fewer channels); the freeze-phase change is not---its
expected gain does not clear seed noise and its reliability cost is
unacceptable for single production runs. For future label rounds the campaign
implies that additional focal-country data does not improve, and likely
harms, the global map; annotation effort should target countries resembling
the current failure profile.

\section{Limitations and planned work}
\label{sec:limitations}
% =====================================================================

Beyond the measurement issues above, four limits are structural.

\emph{Resolution.} At \SI{10}{\metre}, facility \emph{typing} is
information-limited. Pigs-versus-poultry discrimination rests on barn dimensions
and lagoon presence, both marginal at this scale, and the confusion has not
responded to any modelling intervention.

\emph{Temporal collapse.} A single annual median discards phenology, seasonal
lagoon dynamics, and construction change---signals that are plausibly
discriminative and are currently unavailable to the model by construction.

\emph{Cattle.} With 153 training examples, the class is unlearnable and
unmeasurable, and Section~\ref{sec:measurement} shows it is actively harmful to
measurement: it supplies roughly 80\% of the seed variance in the headline
metric. Section~\ref{sec:ablations} shows that a three-class taxonomy costs
nothing measurable. Unless a targeted labelling campaign is undertaken, the class
should be merged away; retaining it as a nominal fourth class with placeholder
metrics is the least defensible option.

\emph{Candidate-stage recall.} Everything downstream is bounded by the
morphometric filters of Table~\ref{tab:morphometry}, which were calibrated on
United States broiler complexes. Facilities that are smaller, differently
proportioned, or built with materials that fragment footprint extraction are
invisible to the entire pipeline, and this failure mode is invisible to our
metrics because such sites never become candidates. We have no estimate of
candidate-stage recall outside the reference region, and consider obtaining one
the most important unmeasured quantity in the system.

The companion document \code{paper/experiments\_justification\_plan.md}
enumerates the experiments required to convert the inherited defaults flagged
with a dagger throughout this paper into justified choices, and records which
have since been run. The measurement upgrades (blind benchmark, spatial-leakage
quantification, seed variance) and the core ablations (band set, spatial context,
optimiser, augmentation, backbone, checkpoint metric, taxonomy) are complete and
reported in Sections~\ref{sec:ablations} and~\ref{sec:measurement}. Three
deployment decisions are settled: a domain-conditional threshold is adopted, and
both model ensembling and geometry fusion are rejected as too small to justify
their cost.

Four items remain open. Three require human annotation and have no compute
substitute: double-annotating evaluation rows to bound the label-noise ceiling,
acquiring \textsc{Cattle} labels (or merging the class, which
Section~\ref{sec:ablations} shows is free), and---most important---measuring
candidate-stage recall outside the reference region. The fourth is a retrain on
spatially blocked splits, for which the split assignment is prepared but the
training run not yet performed.

% =====================================================================
\section{Conclusion}
% =====================================================================

We have described a global pipeline for detecting and typing livestock
facilities from Sentinel-2 imagery, covering \num{157099} candidate clusters
across 167 countries, and reported its experimental record in full. The
methodological findings we would carry to a similar project are: that
sensor-native self-supervised pretraining is worth substantially more than
architectural or optimisation tuning; that class-balanced sampling optimises the
slice sharing the training distribution at the expense of the slice that does
not, and should be evaluated on out-of-domain data before adoption; that
one-sided label campaigns shift decision priors while leaving ranking quality
unchanged, and must be evaluated with threshold-free metrics; and that a silent
identifier-join failure cost us more accuracy than every modelling improvement
we made, which argues for validating spatial joins geometrically rather than by
identifier as a matter of routine.

Two further findings concern measurement rather than modelling, and we would
apply them earliest in any comparable project. Estimate run-to-run variance
before ranking anything: five seeds cost a few GPU-hours here and showed that
most of our prior comparisons sat inside noise. And identify which class drives
that variance before trusting a macro-averaged metric: a class with six held-out
examples supplied roughly 80\% of the movement in the number this project had
been optimising. With those two corrections in place, twenty single-lever
ablations reduce to one change worth making, a third of the input channels shown
to be redundant, and a set of defaults already close to a reasonable optimum.

The system is deployed and its outputs are published, with per-country
performance reported so that users can judge where it is trustworthy. Its
principal remaining weakness is not the model or the benchmark but the stage
before both: we still have no estimate of what fraction of real facilities the
morphometric candidate generator finds outside its reference region, and no
metric we report can see a facility that never became a candidate.

% =====================================================================
\appendix
\section{Reproducibility}
% =====================================================================

\paragraph{Artefacts.} Each release comprises a scored candidate table
(GeoParquet, EPSG:4326 points, carrying per-class probabilities, confidence
tier, split membership, label source, and the imagery tier of
Table~\ref{tab:tiers}), a GeoJSON and CSV export, a metrics summary, and a copy
of the exact training configuration. Releases are registered in a manifest
consumed by an interactive review map.

\paragraph{Software.} PyTorch \citep{paszke2019pytorch} with TorchGeo
\citep{stewart2022torchgeo} for pretrained Sentinel-2 backbones; scikit-learn
\citep{pedregosa2011sklearn} for the Isolation Forest baseline; Google Earth
Engine \citep{gorelick2017gee} for imagery extraction. Experiment tracking uses
MLflow.\footnote{\url{https://mlflow.org}} Candidate tiling uses
H3.\footnote{\url{https://h3geo.org}}

\paragraph{Compute.} Training runs on a single GPU; a full run of the production
configuration completes in under an hour on an NVIDIA RTX 4090 (approximately
55 minutes, 60 seconds per epoch). The ablation campaign of
Section~\ref{sec:ablations} comprised twenty such runs. Full-world inference is
sharded by country across nodes.

\paragraph{Determinism.} Random seeds are fixed at 42 for PyTorch, NumPy, and the
dataset sampler, with per-worker seeding for augmentation. Inference runs in
full precision to keep published probabilities bit-reproducible.

\paragraph{Non-paper data sources.} Registry supervision derives from public
facility registries and from the Farm Transparency
Project\footnote{\url{https://www.farmtransparency.org}} facility database,
supplemented by visual annotation.

\bibliographystyle{plainnat}
\bibliography{references}

\end{document}
